\begin{tabular}{lrr}
\toprule
\textbf{} & \textbf{Co-60} & \textbf{Cs-137} \\
\midrule
archivo & \texttt{Co60\_3M} & \texttt{Cs137\_10M} \\
gammas primarios & 3 M & 10 M \\
eventos con $E_{dep} > 0$ & 852079 & 3531670 \\
fracción que deposita [\%] & 28.4 & 35.3 \\
\midrule
1 scattering Compton [\%] & 75.1 & 73.0 \\
2 scatterings [\%] & 17.1 & 19.3 \\
3 o más [\%] & 5.7 & 7.0 \\
ningún Compton en el centellador [\%] & 2.1 & 0.7 \\
con absorción fotoeléctrica [\%] & 2.9 & 1.3 \\
absorción total ($E_{dep} \approx E_0$) & 87 & 724 \\
\midrule
$\langle L/E_{dep}\rangle$, 50--100 keV & 0.863 & 0.859 \\
$\langle L/E_{dep}\rangle$, $E_{dep} > 700$ keV & 0.954 & -- \\
$\langle L/E_{dep}\rangle$ en la zona del ajuste & 0.954 & 0.935 \\
bordes en luz [keVee] & 921, 1069 & 448 \\
\bottomrule
\end{tabular}
\par\smallskip
\begin{minipage}{0.95\linewidth}\footnotesize Las fracciones por número de scatterings son sobre los eventos con depósito. «Ningún Compton en el centellador» son casi todos eventos cuyo depósito lo hace un electrón Compton nacido en la tapa de aluminio que entra al plástico (mediana de 275 keV en el Cs-137); el resto son fotoeléctricos directos. La zona del ajuste es $E_{dep} > 0{,}7\,E_{\text{ref}}$, que es donde se promedia el quenching que convierte $D_{\text{luz}}$ en $D_{\text{dep}}$.\end{minipage}
